\documentclass{article}
\usepackage{graphicx} % Required for inserting images
\usepackage[portuguese]{babel}
\usepackage{lipsum}
\usepackage{indentfirst}
\usepackage{microtype}
\usepackage[margin=1in]{geometry}
\usepackage{amsmath}
\usepackage{siunitx}
\usepackage{caption}
\usepackage[brazilian]{cleveref}
\usepackage{circuitikz}
\usepackage{pgfplots}
\usepackage{pgfplotstable}
\usepackage{booktabs}
\pgfplotsset{compat=1.18}
\usetikzlibrary{fit, calc}

\addto\captionsportuguese{\renewcommand{\abstractname}{Sumário}} \addto\captionsportuguese{\renewcommand{\contentsname}{Índice}}

\begin{document}
a) Determinar experimentalmente a resistência de entrada de um circuito (por regressão 
linear de um conjunto de pares de valores tensão-corrente) e analiticamente (isto é, por análise 
de circuitos)  
b) Analogamente, determinar a resistência de saída de uma fonte de tensão real 
(simulada)  
c) Determinar o valor duma resistência a partir do seu código de cores
\end{document}
